\section{Mengubah Struktur: ALTER TABLE}

Setelah tabel dibuat, kebutuhan sering berubah --- perlu kolom baru, mengubah tipe data, atau menghapus kolom yang tidak digunakan. Perintah \code{ALTER TABLE} memungkinkan modifikasi struktur tabel tanpa menghapus dan membuat ulang tabel. Ini sangat penting di lingkungan produksi karena data yang sudah ada tetap terjaga. Operasi utama: \code{ADD COLUMN} (menambah kolom), \code{MODIFY COLUMN} (mengubah tipe/constraint), \code{DROP COLUMN} (menghapus kolom), dan \code{RENAME COLUMN} (mengganti nama kolom) \cite{mysqlDocs}.

\begin{webcode}[language=SQL, caption={Contoh penggunaan ALTER TABLE}]
-- Menambah kolom baru
ALTER TABLE produk ADD COLUMN berat DECIMAL(8,2) AFTER stok;

-- Mengubah tipe data kolom
ALTER TABLE produk MODIFY COLUMN nama VARCHAR(255) NOT NULL;

-- Mengganti nama kolom
ALTER TABLE produk RENAME COLUMN dibuat TO created_at;

-- Menghapus kolom
ALTER TABLE produk DROP COLUMN berat;

-- Menambah index
ALTER TABLE produk ADD INDEX idx_kategori (kategori);

-- Menambah foreign key
ALTER TABLE pesanan ADD CONSTRAINT fk_user
  FOREIGN KEY (user_id) REFERENCES pengguna(id);
\end{webcode}

\begin{table}[ht]
\centering
\begin{tabular}{|P{4cm}|P{8.5cm}|}
\hline
\textbf{Perintah ALTER} & \textbf{Fungsi} \\
\hline
\code{ADD COLUMN} & Menambah kolom baru \\
\hline
\code{MODIFY COLUMN} & Mengubah tipe data atau constraint \\
\hline
\code{DROP COLUMN} & Menghapus kolom (dan datanya!) \\
\hline
\code{RENAME COLUMN} & Mengganti nama kolom \\
\hline
\code{ADD INDEX} & Menambah index untuk performa \\
\hline
\code{ADD CONSTRAINT} & Menambah foreign key atau constraint \\
\hline
\end{tabular}
\caption{Ringkasan perintah ALTER TABLE}
\end{table}

